% =====================================================================
%  Szakdolgozat – főfájl
%  Fordítás: latexmk (XeLaTeX + Biber), lásd .latexmkrc
%  VS Code-ban: mentéskor automatikusan fordul (LaTeX Workshop)
% =====================================================================
\documentclass[12pt,a4paper,oneside,openany]{report}

% =====================================================================
%  Formai beállítások a „A szakdolgozat formai és tartalmi követelményei”
%  dokumentum alapján (a pontszámok arra utalnak)
% =====================================================================

% --- Betűtípus (24. pont: Times New Roman) – XeLaTeX kell hozzá -------
\usepackage{fontspec}
\setmainfont{Times New Roman}
\setmonofont{Consolas}[Scale=0.88]   % programkódhoz

% --- Nyelv ------------------------------------------------------------
\usepackage[english,magyar]{babel}   % az utolsó a fő nyelv
\usepackage{csquotes}                 % \enquote{...} -> „...”

% --- Margók (25. pont: bal 3 cm, többi 2,5 cm), A4 --------------------
\usepackage[a4paper,left=3cm,right=2.5cm,top=2.5cm,bottom=2.5cm]{geometry}

% --- Sortávolság (25. pont: másfeles) ---------------------------------
% A 12pt-s report osztályban a lábjegyzet (\footnotesize) eleve 10pt.
\usepackage{setspace}
\onehalfspacing

% --- Címsorok (26. pont: fejezetcím középre, a többi balra) -----------
% (27. pont: minden fejezet új oldalon kezdődik – ezt a report osztály adja)
\usepackage{titlesec}
\titleformat{\chapter}{\normalfont\Large\bfseries\centering}{\thechapter.}{0.6em}{}
\titleformat{name=\chapter,numberless}{\normalfont\Large\bfseries\centering}{}{0pt}{}
\titlespacing*{\chapter}{0pt}{0pt}{1.5em}
\titleformat{\section}{\normalfont\large\bfseries}{\thesection.}{0.6em}{}
\titleformat{\subsection}{\normalfont\normalsize\bfseries}{\thesubsection.}{0.6em}{}
\setcounter{secnumdepth}{2}
\setcounter{tocdepth}{2}

% --- Ábrák, táblázatok (9–10. pont: számozva, címmel) -----------------
\usepackage{graphicx}
\graphicspath{{abrak/}}
\usepackage{booktabs}
\usepackage{tabularx}
\usepackage{array}
\usepackage[font=small,labelfont=bf]{caption}
\usepackage{float}
\usepackage{enumitem}
\setlist{nosep}

% --- Programkód -------------------------------------------------------
\usepackage{xcolor}
\usepackage{listings}
\renewcommand{\lstlistingname}{kódrészlet}
\renewcommand{\lstlistlistingname}{Kódrészletek jegyzéke}
\lstset{
  basicstyle=\ttfamily\small\setstretch{1},
  keywordstyle=\color{blue!60!black}\bfseries,
  commentstyle=\color{green!40!black}\itshape,
  stringstyle=\color{red!50!black},
  numbers=left, numberstyle=\tiny\color{gray}, numbersep=6pt,
  frame=single, rulecolor=\color{gray!50},
  breaklines=true, showstringspaces=false, tabsize=4,
  captionpos=b, columns=fullflexible, keepspaces=true,
}
\lstdefinelanguage{TypeScript}{
  language=Java,
  morekeywords={const,let,async,await,function,interface,type,export,import,from,useState,useEffect},
  morestring=[b]', morestring=[b]`,
}

% --- Irodalomjegyzék (7., 8., 11. pont) -------------------------------
% Egységes, számozott hivatkozás: [3], szó szerinti idézetnél oldalszámmal: [3, 45. o.]
\usepackage[backend=biber,style=numeric,sorting=none,urldate=iso,seconds=true]{biblatex}
\addbibresource{irodalom.bib}
\DefineBibliographyStrings{magyar}{urlseen={Letöltve:}}

% --- Mellékletek, PDF lapok beillesztése (feladatkiírás stb.) ---------
\usepackage{pdfpages}

% --- Kattintható hivatkozások (utolsónak) -----------------------------
\usepackage[hidelinks]{hyperref}

% --- Saját parancsok --------------------------------------------------
% Kitöltendő helyek jelölése: pirosan látszik a PDF-ben, hogy ne maradjon benne
\newcommand{\todo}[1]{\textcolor{red}{\textbf{[#1]}}}
% Ha a megadott PDF létezik, beilleszti; különben egy helyőrző oldalt tesz be
\newcommand{\beillesztPDF}[2]{%
  \IfFileExists{#1}{\includepdf[pages=-]{#1}}{%
    \thispagestyle{empty}\null\vfill
    \begin{center}\Large\textcolor{red}{\textbf{[#2]}}\\[1em]
    \normalsize Ide kerül az aláírt, beszkennelt lap: \texttt{\detokenize{#1}}\end{center}\vfill\clearpage}}

% =====================================================================
%  A dolgozat adatai – ezeket töltsd ki! (a címlap innen dolgozik)
% =====================================================================
\newcommand{\dolgozatCim}{Intelligens online vizsgáztató rendszer fejlesztése automatikus kódértékeléssel és valós idejű felügyelettel}
\newcommand{\hallgatoNev}{\todo{Hallgató neve}}
\newcommand{\hallgatoSzak}{Programtervező informatikus BSc}
\newcommand{\hallgatoEmail}{kisslacika20@gmail.com} % ha egyetemi címet kell megadni, cseréld le
\newcommand{\konzulensNev}{\todo{Konzulens neve}}
\newcommand{\konzulensBeosztas}{\todo{beosztása, pl. egyetemi docens}}
\newcommand{\konzulensIntezet}{\todo{Intézet, pl. Informatikai Intézet}}
\newcommand{\egyetem}{Miskolci Egyetem}
\newcommand{\kar}{Gépészmérnöki és Informatikai Kar}
\newcommand{\hely}{Miskolc}
\newcommand{\ev}{2026}


\begin{document}

% --- 1–3. oldal: aláírt űrlapok (1–3. pont) ---------------------------
% Szkenneld be PDF-be, és tedd a mellekletek/ mappába ezekkel a nevekkel.
\beillesztPDF{mellekletek/feladatkiiras.pdf}{Szakdolgozat feladatkiírás lap}
\beillesztPDF{mellekletek/konzultacios_lap.pdf}{Konzultációs lap}
\beillesztPDF{mellekletek/eredetisegi_nyilatkozat.pdf}{Eredetiségi nyilatkozat}

% --- Belső címlap (4. pont) -------------------------------------------
% Belső címlap: cím, hallgató neve, szakja, e-mail címe, konzulens neve,
% beosztása, intézete, a dolgozat helye (Miskolc) és éve (4. pont)
\begin{titlepage}
\begin{center}
  {\large\textsc{\egyetem}}\\[0.3em]
  {\large \kar}\\[0.3em]
  {\large \konzulensIntezet}

  \vspace*{\stretch{2}}

  {\Large\textbf{SZAKDOLGOZAT}}\\[2em]
  \begin{spacing}{1.2}
  {\LARGE\bfseries \dolgozatCim\par}
  \end{spacing}

  \vspace*{\stretch{3}}

  \begin{tabular}{@{}p{0.45\textwidth}@{\hspace{0.1\textwidth}}p{0.45\textwidth}@{}}
    \textbf{Készítette:}           & \textbf{Konzulens:} \\
    \hallgatoNev                   & \konzulensNev \\
    \hallgatoSzak                  & \konzulensBeosztas \\
    \texttt{\hallgatoEmail}        & \konzulensIntezet \\
  \end{tabular}

  \vspace*{\stretch{1}}

  {\large \hely, \ev}
\end{center}
\end{titlepage}


% --- Tartalomjegyzék (5. pont) ----------------------------------------
\pagenumbering{arabic}
\setcounter{page}{1}
\tableofcontents

% --- (Opcionális) köszönetnyilvánítás a bevezetés előtt (14. pont) ----
% \input{fejezetek/00_koszonet}

% --- Főszöveg ---------------------------------------------------------
\chapter{Bevezetés}
\label{ch:bevezetes}
% Terjedelem: kb. 3–4 oldal (~5 000 karakter)
% 15. pont: a témaválasztás indokai + milyen elméleti és gyakorlati kérdéseket vizsgál a dolgozat.
% 17. pont: összefüggő szöveg, NE felsorolás!

\section{A témaválasztás indoklása}
% - Az online / digitális vizsgáztatás elterjedése (COVID-19 után különösen), a felsőoktatás igényei.
% - Informatikai képzésben a programozási feladatok kézi javítása időigényes és nem egységes.
% - Meglévő rendszerek (Moodle Quiz, CodeRunner stb.) hiányosságai: kódfuttatás biztonsága,
%   valós idejű felügyelet hiánya, összetett beállítás.
% - Saját motiváció (saját vélemény is kell – 19. pont).
\todo{Szöveg}

\section{Célkitűzések}
% Egy olyan webes prototípus elkészítése, amely:
%   vegyes feladattípusokat kezel (kód, egy/több választós, igaz–hamis, rövid válasz, esszé),
%   a programozási feladatokat biztonságos Docker-sandboxban, tesztesetekkel automatikusan pontozza,
%   valós időben jelzi az oktatónak, ha a hallgató elhagyja a vizsgaablakot,
%   oktatói felületet ad a kézi javításhoz és az eredmények exportálásához.
% Vizsgált kérdések: hogyan lehet nem megbízható kódot biztonságosan futtatni?
% Milyen pontozási szabályok igazságosak? Mennyit ér a böngészőalapú felügyelet?
\todo{Szöveg}

\section{A dolgozat felépítése}
% Röviden, fejezetenként (\ref{ch:szakirodalom}. fejezet ..., \ref{ch:technologiak}. fejezet ...).
\todo{Szöveg}

\chapter{Az online vizsgáztatás szakirodalmi háttere}
\label{ch:szakirodalom}
% Terjedelem: kb. 10–12 oldal (~18 000 karakter)
% 13. pont: a dolgozat 50–60%-a szakirodalomra épüljön -> ez és a 3. fejezet adja a gerincét.
% Minden állításnál hivatkozás: \cite{kulcs}, szó szerinti idézetnél: \cite[12.~o.]{kulcs}

\section{Számítógéppel támogatott értékelés}
% E-learning, LMS-ek, a számítógépes értékelés (CAA) fogalma, előnyei/hátrányai.
% Formatív vs. szummatív értékelés. Források: \cite{bull2004blueprint}
\todo{Szöveg}

\section{Feladattípusok és automatikus pontozásuk}
% Zárt (választós, igaz–hamis), félig nyílt (rövid válasz), nyílt (esszé) feladatok.
% Több jó válaszos feladatok pontozási módszerei (részpontozás, levonás a rossz jelölésért).
\todo{Szöveg}

\section{Programozási feladatok automatikus értékelése}
% Teszteset-alapú értékelés, nyilvános és rejtett tesztek, ismert rendszerek (pl. CodeRunner,
% DOMjudge, Web-CAT). Források: \cite{douce2005automatic,ihantola2010review,paiva2022automated}
\todo{Szöveg}

\section{Online felügyelet (proctoring) és akadémiai integritás}
% Proctoring típusai (élő, rögzített, automatizált), böngészőalapú megoldások,
% adatvédelmi és etikai kérdések. Források: \cite{atoum2017automated,nigam2021systematic,coghlan2021proctor}
\todo{Szöveg}

\section{Meglévő rendszerek összehasonlítása}
% Táblázatos összevetés + szöveges értékelés (saját vélemény!). A táblázatra hivatkozni kell:
A meglévő rendszerek főbb jellemzőit a \ref{tab:osszehasonlitas}.~táblázat foglalja össze.

\begin{table}[H]
  \centering
  \caption{Online vizsgáztató rendszerek összehasonlítása \todo{ellenőrizendő}}
  \label{tab:osszehasonlitas}
  \begin{tabularx}{\textwidth}{@{}l>{\centering\arraybackslash}X>{\centering\arraybackslash}X>{\centering\arraybackslash}X>{\centering\arraybackslash}X@{}}
    \toprule
    Szempont                  & Moodle Quiz & CodeRunner & \todo{?} & Saját rendszer \\
    \midrule
    Vegyes feladattípusok     & igen        & részben    &          & igen \\
    Kódfuttatás sandboxban    & nem         & igen       &          & igen (Docker) \\
    Rejtett tesztesetek       & --          & igen       &          & igen \\
    Valós idejű felügyelet    & nem         & nem        &          & igen \\
    \bottomrule
  \end{tabularx}
\end{table}

\chapter{Alkalmazott technológiák}
\label{ch:technologiak}
% Terjedelem: kb. 8–10 oldal (~14 000 karakter)
% Nem csak felsorolás: mindegyiknél MIÉRT ezt választottad, mi az alternatíva.

\section{Webes architektúra és a REST elv}
% Kliens–szerver, REST megszorításai \cite{fielding2000architectural}, HTTP metódusok.
\todo{Szöveg}

\section{A backend technológiái}
\subsection{FastAPI és az ASGI}
% Típusannotációk, Pydantic validáció, automatikus OpenAPI/Swagger \cite{fastapi_docs}.
\todo{Szöveg}

\subsection{Adatbázis: PostgreSQL és SQLAlchemy ORM}
% Relációs modell, ORM előnyei \cite{sqlalchemy_docs}.
\todo{Szöveg}

\subsection{Hitelesítés: JWT és bcrypt}
% Állapotmentes token \cite{rfc7519}, jelszó-hash sózással \cite{provos1999bcrypt}.
\todo{Szöveg}

\section{Valós idejű kommunikáció: WebSocket és Socket.IO}
% \cite{rfc6455,socketio_docs}
\todo{Szöveg}

\section{A frontend technológiái: React, Next.js, Tailwind CSS}
% Komponensalapú UI, App Router, kliensoldali állapotkezelés \cite{nextjs_docs}.
\todo{Szöveg}

\section{Konténerizáció és sandboxing: Docker}
% Konténer vs. virtuális gép, névterek (namespaces), cgroups, képességek (capabilities)
% \cite{merkel2014docker,docker_docs}. Nem megbízható kód futtatásának kockázatai \cite{owasp_top10}.
\todo{Szöveg}

\chapter{A rendszer tervezése}
\label{ch:tervezes}
% Terjedelem: kb. 7–8 oldal (~12 000 karakter)

\section{Követelmények}
\subsection{Funkcionális követelmények}
% Szerepkörök: hallgató, oktató (meghívókóddal regisztrál).
% Oktató: vizsga létrehozása (időablak + időtartam), feladatok típusonként, élő monitor,
%   eredmények, kézi pontozás + visszajelzés, CSV export, újraértékelés.
% Hallgató: vizsga indítása, visszaszámláló, automatikus mentés, próbafuttatás, beadás, eredmény.
\todo{Szöveg}

\subsection{Nem funkcionális követelmények}
% Biztonság (megoldókulcs sosem jut a klienshez, sandbox), használhatóság, hordozhatóság (Docker).
\todo{Szöveg}

\section{Architektúra}
% Háromrétegű felépítés: Next.js kliens – FastAPI + Socket.IO szerver – PostgreSQL,
% mellette a Docker sandbox. Ide jön egy ábra (abrak/architektura.png):
\todo{Szöveg, ábra: \ref{fig:architektura}.~ábra}

\begin{figure}[H]
  \centering
  % \includegraphics[width=0.9\textwidth]{architektura}
  \fbox{\parbox[c][4cm][c]{0.8\textwidth}{\centering Helyőrző: architektúra ábra}}
  \caption{A rendszer architektúrája}
  \label{fig:architektura}
\end{figure}

\section{Adatmodell}
% Entitások: User, Exam, Question, Attempt, Answer – ER-diagram + magyarázat.
% Fontos tervezési döntések: egy hallgató/vizsga = egy Attempt (UNIQUE), egy válasz/feladat (UNIQUE),
% answer_key JSON típusonként eltérő szerkezettel, manual_score felülírja az auto_score-t.
A rendszer adattábláit a \ref{tab:adatmodell}.~táblázat foglalja össze.

\begin{table}[H]
  \centering
  \caption{Az adatmodell táblái}
  \label{tab:adatmodell}
  \begin{tabularx}{\textwidth}{@{}lX@{}}
    \toprule
    Tábla & Szerep és fontosabb mezők \\
    \midrule
    \texttt{users}     & felhasználók; e-mail, név, jelszó-hash, szerepkör (\texttt{student}/\texttt{teacher}) \\
    \texttt{exams}     & vizsgák; cím, időtartam (perc), kezdési és záró időpont, létrehozó oktató \\
    \texttt{questions} & feladatok; típus, pontszám, válaszlehetőségek, megoldókulcs (JSON) \\
    \texttt{attempts}  & kitöltések; kezdés, határidő, befejezés, szabálysértések száma \\
    \texttt{answers}   & válaszok; válasz (JSON), automatikus és kézi pontszám, visszajelzés, állapot \\
    \bottomrule
  \end{tabularx}
\end{table}

\section{A REST API terve}
\begin{table}[H]
  \centering
  \caption{A REST API fontosabb végpontjai}
  \label{tab:api}
  \small
  \begin{tabularx}{\textwidth}{@{}lp{6.2cm}X@{}}
    \toprule
    Metódus & Végpont & Funkció \\
    \midrule
    POST   & \texttt{/auth/register}, \texttt{/auth/login} & regisztráció, bejelentkezés (JWT) \\
    POST   & \texttt{/exams/} & vizsga létrehozása (oktató) \\
    POST   & \texttt{/exams/\{id\}/questions} & feladat hozzáadása \\
    POST   & \texttt{/exams/\{id\}/attempt} & kitöltés indítása (hallgató) \\
    PUT    & \texttt{/attempts/\{id\}/answers/\{qid\}} & válasz mentése \\
    POST   & \texttt{/attempts/\{id\}/run} & kód próbafuttatása \\
    POST   & \texttt{/attempts/\{id\}/finish} & végleges beadás \\
    GET    & \texttt{/exams/\{id\}/results(.csv)} & eredmények, CSV export \\
    PATCH  & \texttt{/answers/\{id\}/grade} & kézi pontozás \\
    POST   & \texttt{/exams/\{id\}/regrade} & újraértékelés \\
    \bottomrule
  \end{tabularx}
\end{table}

\section{Biztonsági terv}
% Fenyegetések és ellenintézkedések: jogosultság (csak a saját vizsga), megoldókulcs elrejtése,
% Socket.IO JWT-hitelesítés, sandbox korlátok.
\todo{Szöveg}

\chapter{Megvalósítás}
\label{ch:megvalositas}
% Terjedelem: kb. 9–10 oldal (~15 000 karakter)
% Kódrészletből csak a lényeges, rövid részeket! A teljes kód mellékletbe / repóba.

\section{A backend felépítése}
% Mappaszerkezet: api/ (auth, exams, attempts, results), core/ (database, security, timeutils),
% models/, schemas/ (Pydantic), services/ (grading, sandbox, socket_server).
\todo{Szöveg}

\section{Hitelesítés és jogosultságkezelés}
% bcrypt hash, JWT (HS256, 1 nap), oktatói meghívókód, szerepkör-ellenőrzés.
\todo{Szöveg}

\section{A vizsgafolyamat és az időkezelés}
% Kezdés -> határidő = min(kezdés + időtartam, vizsga záró időpontja); 30 mp-es türelmi idő;
% automatikus mentés; lejáratkor automatikus lezárás (finalize_if_expired).
\todo{Szöveg}

\section{Az automatikus pontozás}
% Típusonként. A több jó válaszos feladat képlete:
A több helyes választ is tartalmazó feladatoknál a rendszer részpontot ad, a hibás jelöléseket
pedig levonja:
\begin{equation}
  P = \max\!\left(0,\ \frac{|S \cap H| - |S \setminus H|}{|H|}\right) \cdot P_{\max},
  \label{eq:tobbvalasz}
\end{equation}
ahol $S$ a megjelölt, $H$ a helyes válaszok halmaza, $P_{\max}$ pedig a feladat pontértéke.
% Rövid válasz: kis/nagybetű- és szóközfüggetlen, számoknál 3,14 = 3.14.
% Kódfeladat: pont = sikeres tesztek / összes teszt * pont.
\todo{Szöveg}

\section{Biztonságos kódfuttatás Docker-sandboxban}
% Egy konténer / kiértékelés, runner script, nobody felhasználó, kimenet- és időkorlát.
A konténerre alkalmazott korlátozásokat a \ref{tab:sandbox}.~táblázat, a konténer indítását
a \ref{lst:sandbox}.~kódrészlet mutatja be.

\begin{table}[H]
  \centering
  \caption{A sandbox konténer korlátozásai}
  \label{tab:sandbox}
  \begin{tabularx}{\textwidth}{@{}lX@{}}
    \toprule
    Korlátozás & Érték, cél \\
    \midrule
    Hálózat             & nincs (\texttt{network\_mode="none"}) – adatszivárgás ellen \\
    Memória             & 128 MB \\
    CPU                 & 0,5 mag \\
    Folyamatok száma    & legfeljebb 64 (fork bomb ellen) \\
    Fájlrendszer        & csak olvasható; \texttt{/tmp}: 16 MB-os tmpfs \\
    Jogosultságok       & minden capability eldobva, \texttt{no-new-privileges}; a kód \texttt{nobody}-ként fut \\
    Időkorlát           & 5 mp tesztesetenként \\
    Kimenet             & legfeljebb 64 KB tesztesetenként \\
    \bottomrule
  \end{tabularx}
\end{table}

\begin{lstlisting}[language=Python,caption={A sandbox konténer indítási paraméterei (\texttt{sandbox.py}, részlet)},label={lst:sandbox}]
run_kwargs = dict(
    image=SANDBOX_IMAGE,
    command=["python", "-c", _RUNNER],
    network_mode="none",
    mem_limit="128m",
    nano_cpus=500_000_000,
    pids_limit=64,
    read_only=True,
    tmpfs={"/tmp": "size=16m,mode=1777"},
    cap_drop=["ALL"],
    cap_add=["SETUID", "SETGID", "KILL"],
    security_opt=["no-new-privileges"],
)
\end{lstlisting}

\section{Valós idejű felügyelet}
% Kliens: useProctoring hook (blur + visibilitychange, dupla számolás kiszűrése).
% Szerver: Socket.IO események (connect JWT-vel, join_exam, report_violation, exam_finished),
% szobák vizsgánként, a szabálysértések az adatbázisban tárolódnak.
\todo{Szöveg}

\section{A felhasználói felület}
% Képernyőképek (abrak/ mappába): bejelentkezés, hallgatói műszerfal, vizsgaoldal kódszerkesztővel,
% oktatói feladatszerkesztő, élő monitor, eredmények. Minden képre hivatkozz a szövegben!
\todo{Szöveg és képernyőképek}

\chapter{Tesztelés és értékelés}
\label{ch:teszteles}
% Terjedelem: kb. 4–5 oldal (~6 000 karakter)

\section{Funkcionális tesztelés}
% Tesztforgatókönyvek táblázatban: lépés – elvárt eredmény – tényleges eredmény.
% Pl. vizsga lejárata utáni beadás, második kitöltés tiltása, más oktató vizsgájának elérése.
\todo{Szöveg}

\section{A sandbox biztonsági tesztelése}
% Rosszindulatú kódok és a rendszer reakciója: végtelen ciklus, fork bomb, nagy memóriafoglalás,
% hálózati kapcsolat, fájlírás, óriási kimenet. Ide jó egy táblázat.
\todo{Szöveg}

\section{Korlátok}
% Saját értékelés: a böngészőalapú proctoring könnyen kijátszható (második eszköz),
% csak Python támogatott, nincs terheléses teszt, stb.
\todo{Szöveg}

\chapter{Összefoglalás}
\label{ch:osszefoglalas}
% Terjedelem: kb. 2–3 oldal (~4 000 karakter)
% 21. pont: elméleti és gyakorlati tanulságok, következtetések, eredmények összegzése.
% 19. pont: saját vélemény!

\section{Eredmények}
\todo{Szöveg}

\section{Továbbfejlesztési lehetőségek}
% Több programozási nyelv, plágiumkeresés, webkamerás proctoring, LMS-integráció (LTI),
% skálázás (feladatsor/worker a kódfuttatáshoz), statisztikák.
\todo{Szöveg}


% --- Idegen nyelvű összefoglaló, 1–2 oldal (22. pont) -----------------
% Idegen nyelvű összefoglaló, 1–2 oldal (22. pont)
% Számozatlan fejezet, de bekerül a tartalomjegyzékbe.
\chapter*{Summary}
\addcontentsline{toc}{chapter}{Summary}
\begin{otherlanguage}{english}
\todo{English summary, 1--2 pages: motivation, goals, the developed system, results, future work.}
\end{otherlanguage}


% --- Irodalomjegyzék és mellékletek a végén (23. pont) ----------------
\printbibliography[heading=bibintoc,title={Irodalomjegyzék}]

\listoffigures
\addcontentsline{toc}{chapter}{Ábrajegyzék}
\listoftables
\addcontentsline{toc}{chapter}{Táblázatjegyzék}

\appendix
% Mellékletek (23. pont: a dolgozat végén). \appendix után a fejezetek betűjelet kapnak.
\chapter{Telepítési és futtatási útmutató}
\label{app:telepites}
% A README alapján: Docker, Python, Node.js; 3 terminál; környezeti változók.
\todo{Szöveg}

\chapter{A forráskód elérhetősége}
% Pl. GitHub link vagy a mellékelt adathordozó tartalma.
\todo{Szöveg}


\end{document}
